% VitalStream, the waveform viewer, in full: installing, running, the page, opening records and
% folders, the lanes, time, the settings, deploying, adapters, testing.
% Sources: viewers/vitalstream/README.md, viewers/vitalstream/docs/*.md, docs/wfdb-collection.md,
% docs/edf.md,
% checked against viewers/vitalstream/web/index.html, web/js/app.js and render.js,
% rust/src/main.rs, settings.rs, server.rs and auth.rs, install.sh and install.ps1, and the
% folder limits of libs/engine/src/sources/wfdb_folder.rs.
\chapter{VitalStream}
\label{chap:vitalstream}

This chapter describes VitalStream, the waveform viewer, as a program on its own: installing
it, running it, its page, and its settings. It is written for researchers who run the viewer on
their own computer or a server, and for administrators.

VitalStream shows waveform recordings in a browser: every channel, the vital-sign trends, the
alarms or annotations, and the breaks in the data, on one time axis, from weeks down to single
samples. Long, high-frequency recordings are its strength. ICU bedside monitoring, Holter ECGs,
EEGs and sleep studies of hours, days or weeks open at once and stay fluid while zooming.
VitalStream is the waveform viewer of Medorama, which opens recordings in it at the moment chosen
on a patient's page (\cref{chap:viewers-in-medorama}), and it is also a complete program on its
own. Its page has a light and a dark theme, and follows the system's setting until one is
chosen.

VitalStream adds nothing that is not in the recording. Samples are never resampled, interpolated
or stitched. Where the recording is ambiguous (a break's length, a record without a start time),
the page says so. Times are kept in UTC; the time zone in which they are shown is a setting of
the page only.

\screenshot{vitalstream.png}{VitalStream in the dark theme: a synthetic ICU recording (a WFDB collection), with the ECG, arterial pressure, pleth and respiration waveforms, an alarm, the vital-sign trends and the values at the cursor.}{fig:vitalstream-main}

% ==========================================================================================
\section{What it reads}
\label{sec:vitalstream-reads}

\begin{table}[htbp]
\centering
\caption{The recordings VitalStream reads.}
\label{tab:vitalstream-reads}
\begin{tabularx}{\linewidth}{@{}LP{4.6cm}@{}}
\toprule
\tabhead Recordings & \tabhead How \\
\midrule
One WFDB record (PhysioNet's format): its \file{.hea} header file, single- or multi-segment &
  Read by the viewer itself \\
A folder of WFDB records, its subfolders included, read as \emph{one session} on one time axis;
  a WFDB collection (records in time order with their vital signs, alarms and breaks) & Read by
  the viewer itself \\
An EDF, EDF+, BDF or BDF+ file (\file{.edf}, \file{.bdf}): sleep studies, EEG, long-term ECG &
  Read by the viewer itself (\cref{sec:vitalstream-edf}) \\
Other formats, licensed ones among them (\file{.Stp} files of BedMasterEx, through the
  stp-adapter add-on, which is private and available to institutions under agreement) &
  Through an installed \emph{adapter} (\cref{sec:vitalstream-adapters}) \\
\bottomrule
\end{tabularx}
\end{table}

WFDB collections, multi-segment records and the files a collection holds (\file{index.csv},
\file{gaps.csv}, \file{vitals.csv}, \file{events.csv}, the \file{.mrk} markers) are described in
\cref{chap:formats-reference}. VitalStream assumes nothing that a record does not state. A
record whose header gives no sampling frequency is not read: opened alone, it is refused, and in
a folder of records it is left out and counted in the notes. A signal with a gain of 0 is shown
in stored counts, with the unit \code{adu}, and a signal whose header gives no units has none.
WFDB's old defaults of 250\,Hz, a gain of 200 and millivolts are not used.

% ==========================================================================================
\section{EDF, EDF+, BDF and BDF+ files}
\label{sec:vitalstream-edf}

The three viewers read files of the European Data Format family directly, without an add-on:
EDF and BDF (16-bit and 24-bit samples), and EDF+ and BDF+, which add annotations and
interrupted recordings (\cref{sec:recordings-edf}). An EDF file is one recording. A file is
opened as EDF when its name ends in \file{.edf} or \file{.bdf}, in any letter case; choose
\ui{Open record\dots} and pick the file, or type its path in the dialog. The built-in reader takes
precedence over any adapter that lists these extensions. \Cref{tab:vitalstream-edf} lists what
the viewer shows; the exact reading rules are specified in \cref{sec:formats-reference-edf}.

\begin{xltabular}{\linewidth}{@{}P{5.0cm}L@{}}
\caption{An EDF file in VitalStream.}
\label{tab:vitalstream-edf}\\
\toprule
\tabhead In the file & \tabhead In the viewer \\
\midrule
\endfirsthead
\toprule
\tabhead In the file & \tabhead In the viewer \\
\midrule
\endhead
\bottomrule
\endfoot
Each ordinary signal & A channel, with the signal's label and physical dimension (\code{ECG},
  \code{mV}). \\
Samples per data record divided by the duration of a data record & The channel's sampling rate.
  Signals may have different rates. \\
Digital values & Physical values, scaled with the signal's physical and digital minimum and
  maximum. \\
The start date and time in the header & The time axis. The time is taken as UTC, since EDF does
  not state a time zone. \\
EDF+ and BDF+ annotations & Events in the \ui{Annotations} lane, with their text and duration.
  Hover over one to read it. \\
An interrupted recording (EDF+D) & One stretch of data per uninterrupted run, with a break
  between runs (\code{not recorded}). \\
A second signal with the same label & A channel named with its count: \code{ECG (2)}. \\
\end{xltabular}

The \ui{Recording} panel shows the format as \code{EDF}, \code{EDF+}, \code{BDF} or
\code{BDF+}, and the notes that apply, for example when a file ends partway through a data
record. The header's patient and recording identification fields are never read into the index
and never shown. Under Medorama, in Privacy Mode, the annotations' texts are treated like the text
of notes (\cref{sec:protecting-edf}).

\screenshot{vitalstream-sleep.png}{VitalStream on the demo's EDF+ sleep study, with its
EEG, EOG, ECG, airflow, effort and SpO$_2$ channels and its Annotations lane (dark
theme).}{fig:vitalstream-sleep}

Two limits apply. A folder of EDF files is not yet opened as one session: \ui{Open multiple
records\dots} reads the WFDB records of a folder, so open EDF files one at a time. An EDF+ file
that holds only annotations (some sleep-scoring files do) has no signal to draw and is
refused.

% ==========================================================================================
\section{Three viewers, one specification}
\label{sec:vitalstream-three}

VitalStream has three implementations: the same page, the same requests and answers, the same
index, in three languages. Use whichever fits the machine; several can run side by side.

\begin{table}[htbp]
\centering
\caption{The three VitalStream viewers.}
\label{tab:vitalstream-three}
\begin{tabularx}{\linewidth}{@{}P{3.2cm}LLL@{}}
\toprule
& \tabhead Python & \tabhead Rust & \tabhead C++ \\
\midrule
Command & \cmd{vitalstream-py} & \cmd{vitalstream-rs} & \cmd{vitalstream-cpp} \\
Port (default) & 8050 & 8060 & 8070 \\
Folder & \file{python/vitalstream} & \file{rust/} & \file{cpp/} \\
Needs to run & Python 3.9 or later & Nothing (one file) & Nothing but the C++ runtime (one
  file) \\
Indexing a day of ICU data & About a minute & A few seconds & A few seconds \\
Package, library & The Python package \code{vitalstream\_py} & The crate \code{vitalstream-rs},
  with the engine \code{recording-engine} (\file{libs/engine}) & The CMake targets
  \code{vs\_core}, \code{vs\_engine}, \code{vs\_viewer} \\
\bottomrule
\end{tabularx}
\end{table}

The folders are those of \file{viewers/vitalstream}. The three serve the same page
(\file{viewers/vitalstream/web}: the Rust and C++ viewers have it built in, the Python viewer
serves it from its package), answer the same HTTP API with the same JSON
(\cref{chap:viewer-api}), write the same index byte for byte (\cref{sec:vitalstream-index}), and
share one settings file, one password, one index folder and one adapter folder. Every name a
user meets ends in the implementation's suffix, \code{-py}, \code{-rs} or \code{-cpp}; without a
suffix, a name refers to what the three have in common (``the viewer'', \file{vitalstream.toml}).
The three read WFDB records with wfdb-x, in each language.

% ==========================================================================================
\section{Installing}
\label{sec:vitalstream-install}

VitalStream alone is installed from a clone of the repository, with no administrator rights
beyond the compiler, by \file{viewers/vitalstream/install.sh} (Linux, macOS) or
\file{viewers/vitalstream/install.ps1} (Windows); the commands for each operating system are
also at \url{https://medorama.org/docs/install.html}. The script builds one viewer and installs it:

\begin{lstlisting}[style=shell]
git clone https://github.com/alphanumericslab/medorama.git && cd medorama
viewers/vitalstream/install.sh                         # the C++ viewer, into ~/.local/bin
viewers/vitalstream/install.sh --viewer rs             # or the Rust viewer (or: py)
viewers/vitalstream/install.sh --root /data/icu        # it opens nothing outside /data/icu
viewers/vitalstream/install.sh --no-password           # a viewer without password support
viewers/vitalstream/install.sh --no-wfdb               # a viewer without WFDB support (no wfdb-x)
\end{lstlisting}

On Windows, the following command runs the installer in PowerShell (for the C++ viewer, in a
``Developer PowerShell for VS''):

\begin{lstlisting}[style=shell]
powershell -ExecutionPolicy Bypass -File viewers\vitalstream\install.ps1 -Viewer rs -Root D:\data
\end{lstlisting}

\begin{xltabular}{\linewidth}{@{}P{2.9cm}P{1.8cm}P{2.3cm}L@{}}
\caption{The options of VitalStream's installers.}
\label{tab:vitalstream-install}\\
\toprule
{\tabhead install.sh} & {\tabhead install.ps1} & {\tabhead Default} & {\tabhead Effect} \\
\midrule
\endfirsthead
\toprule
{\tabhead install.sh} & {\tabhead install.ps1} & {\tabhead Default} & {\tabhead Effect} \\
\midrule
\endhead
\bottomrule
\endfoot
\opt{-{}-viewer cpp}, \opt{rs} or \opt{py} & \opt{-Viewer} & \code{cpp} & The viewer to build: C++
  (a C++20 compiler, CMake and Ninja; \file{install.sh} installs CMake and Ninja into a build
  environment, \file{viewers/vitalstream/cpp/build/tools}, when they are missing), Rust
  (\file{install.sh} installs Rust in your account when \cmd{cargo} is missing), Python \\
\opt{-{}-prefix DIR} & \opt{-Prefix} & \file{\textasciitilde/.local} (Windows:
  \code{\%LOCALAPPDATA\%\textbackslash vitalstream}) & The program goes into \file{DIR/bin}; the
  Python viewer's environment into \file{DIR/share/vitalstream/venv} (Windows:
  \code{DIR\textbackslash venv}), with a link to its program in \file{DIR/bin} \\
\opt{-{}-root DIR} & \opt{-Root} & \code{\textasciitilde} & The viewer's root, fixed in the program:
  its page never lists or opens anything outside \file{DIR}. \code{\textasciitilde} is the home
  folder of whoever runs the viewer; \code{/} is the whole file system \\
\opt{-{}-no-password} & \opt{-NoPassword} & Password support & Build it without password
  support: no password can be set (\opt{-{}-set-password} and \opt{-{}-clear-password} fail), a
  password already set in the shared file is ignored (the page never asks for one), and it
  refuses to run where one is required \\
\opt{-{}-no-wfdb} & \opt{-NoWfdb} & WFDB support & Build it without WFDB support: it contains no
  wfdb-x code, refuses a WFDB record or folder with a message (``this build has no WFDB
  support''), and reads every other format (\cref{sec:installation-no-wfdb}) \\
\opt{-{}-no-rustup} & & & Fail, rather than install Rust, when \cmd{cargo} is missing (the Rust
  viewer); \file{install.ps1} never installs Rust \\
\opt{-h}, \opt{-{}-help} & & & Print the options and exit \\
\end{xltabular}

The script ends with the program's path and its root, for instance
\code{Installed \textasciitilde/.local/bin/vitalstream-cpp (root: \textasciitilde)}, followed by
\code{; no password support} or \code{; no WFDB support} for a build without them, and says
how to start it. When
\file{DIR/bin} is not on the \code{PATH}, it says so: add it, or call the program by its path.
Run the script again to change the root or the password support.

The installer's \opt{-{}-root} is not the viewer's \opt{-{}-root}: the first is fixed when the
viewer is built and cannot be widened afterwards; the second (\cref{sec:vitalstream-run}) narrows
it each time the viewer starts (\cref{chap:viewers-alone}). The standalone builds go to folders of
their own (\file{target/standalone}, \file{viewers/vitalstream/cpp/build/standalone}), so that
Medorama's viewer in the same clone keeps its own. A build by hand sets the same three options:

\begin{table}[htbp]
\centering
\caption{The build options of VitalStream, by hand.}
\label{tab:vitalstream-build-options}
\begin{tabularx}{\linewidth}{@{}P{2.3cm}LLL@{}}
\toprule
\tabhead Viewer & \tabhead Root & \tabhead Without password support & \tabhead Without WFDB support \\
\midrule
\cmd{vitalstream-rs} & \code{VITALSTREAM\_ROOT=DIR} when it is compiled &
  \opt{-{}-no-default-features -{}-features wfdb} (the Cargo feature \code{password} left out) &
  \opt{-{}-no-default-features -{}-features password} (the Cargo feature \code{wfdb} left out) \\
\cmd{vitalstream-cpp} & CMake \code{-DVS\_ROOT=DIR} & \code{-DVS\_PASSWORD=OFF} &
  \code{-DVS\_WFDB=OFF} (the preset \code{nowfdb}) \\
\cmd{vitalstream-py} & \code{VITALSTREAM\_ROOT=DIR} when the package is installed &
  \code{VITALSTREAM\_PASSWORD=0} & Installed without its \code{[wfdb]} extra \\
\bottomrule
\end{tabularx}
\end{table}

Without them (a developer's \cmd{cargo build}, \cmd{cmake -{}-preset release},
\cmd{pip install}), the root is the whole file system and the password is supported. The
compilers, CMake presets, builds without administrator rights and without internet access are
described in \cref{chap:installation}.

\paragraph{With Medorama.} \cmd{tools/install-medorama.sh} builds Medorama and one viewer
together. That viewer has the whole file system as its root and password support, since Medorama
applies its own rules: it starts the viewer itself, with its own settings, which require a
password (\cref{chap:viewers-in-medorama,chap:administration}).

\paragraph{Update and uninstall.} To update, \cmd{git pull} and run the installer again with the
same options, then restart the viewer. To uninstall, delete \file{\textasciitilde/.local/bin/vitalstream-*}
(and, for the Python viewer, \file{\textasciitilde/.local/share/vitalstream}); the index folder is
\file{\textasciitilde/.cache/vitalstream}, the settings and the password are in
\file{\textasciitilde/.config/vitalstream/}. Recordings are never touched.

% ==========================================================================================
\section{Running}
\label{sec:vitalstream-run}

The following commands start the C++ viewer in the most common ways; the other two viewers take
the same options.
\begin{lstlisting}[style=shell]
vitalstream-cpp                       # opens http://127.0.0.1:8070 in your browser
vitalstream-cpp --set-password        # once, on a shared machine (one password for the viewers)
vitalstream-cpp --no-browser          # on a server; then an SSH tunnel to its port
vitalstream-cpp --root ~/icu          # open nothing outside ~/icu (inside the viewer's root)
vitalstream-cpp --index RECORDING     # index ahead of time
vitalstream-cpp --help                # every option
\end{lstlisting}

The three viewers take the same options. Each option overrides the setting of the same name
(\cref{sec:vitalstream-settings}).

\begin{lstlisting}[style=signature]
vitalstream-cpp [--config FILE] [--host HOST] [--port N] [--timezone ZONE] [--root DIR]...
                [--cache-dir DIR] [--adapters-dir DIR] [--qc-verbose] [--no-browser]
                [--set-password | --clear-password] [--no-auth] [--index PATH] [-V]
\end{lstlisting}

\begin{params}
\item[--config FILE] The settings file. Default: the first that exists of
  \file{viewers/vitalstream/vitalstream.toml} (only for a viewer run from the build folder of a
  clone) and \file{\textasciitilde/.config/vitalstream/vitalstream.toml}; a
  \file{vitalstream.toml} in the working folder is never read unless \opt{-{}-config} names it.
  A \file{FILE} that does not exist is an error (\code{FILE: no such settings file}).
\item[--host HOST] The address to listen on. Default: \code{127.0.0.1}, this machine only.
\item[--port N] The port. Default: 8050 (\cmd{vitalstream-py}), 8060 (\cmd{vitalstream-rs}),
  8070 (\cmd{vitalstream-cpp}).
\item[--timezone ZONE] The display time zone, an IANA name such as \code{America/New\_York}.
  An unknown zone is an error (\code{unknown time zone}).
\item[--root DIR] A folder the viewer may open, inside its root; repeatable. Default: its root
  (the setting \code{data\_roots}).
\item[--cache-dir DIR] The index folder (\cref{sec:vitalstream-index}).
\item[--adapters-dir DIR] The adapter folder (\cref{sec:vitalstream-adapters}).
\item[--qc-verbose] Also mark the data a reader repaired, such as data blocks put back in order
  (\cref{sec:vitalstream-qc}).
\item[--no-browser] Do not open a browser window.
\item[--set-password] Set (or change) the password the page asks for, then exit. It asks for the
  password twice (\code{New VitalStream password:}, \code{Repeat it:}); at least 8 characters.
  A viewer built without password support refuses (\code{this viewer was built without
  password support}).
\item[--clear-password] Remove the password, then exit; refused, as \opt{-{}-set-password}, by a
  viewer built without password support.
\item[--no-auth] Run without asking for the password, even when one is set.
\item[--index PATH] Index \file{PATH} (a record's \file{.hea}, a folder, a file an adapter reads)
  into the index folder, print \code{indexed PATH -> FOLDER/KEY/}, and exit.
\item[-V, --version] Print the program's name and version (\code{vitalstream-py 0.1.0}), then
  exit.
\end{params}

When it starts, the viewer prints the settings file it read, its address, its root, its data
roots, the index folder and the time zone, and whether a password protects it:

\begin{lstlisting}[style=text]
vitalstream-rs: settings from /home/user/.config/vitalstream/vitalstream.toml
VitalStream (Rust): open http://127.0.0.1:8060/
  root: /; data roots: /data/icu; cache: /home/user/.cache/vitalstream; time zone UTC
  no password: anyone who can connect to this port can read your files through it.
  On a shared machine, set one once with: vitalstream-rs --set-password
\end{lstlisting}

Without a settings file it says \code{no settings file (defaults)}. An IPv6 \opt{-{}-host} is
written with brackets in the address (\code{http://[::1]:8050/}). A data root outside the
viewer's root is left out, with a warning; with none left, the viewer
does not start. With the setting \code{require\_password} and no password set (or with
\opt{-{}-no-auth}), it does not start either (\code{a password is required here
(require\_password); set one with -{}-set-password}; for a viewer built without password
support, \code{a password is required here (require\_password), and this viewer was built
without password support}). A viewer built without password support ignores a password already
set, and says at start \code{built without password support: anyone who can connect to this
port can read the files under its root}. The viewer stops with \key{Ctrl}+\key{C}.

% ==========================================================================================
\section{The page}
\label{sec:vitalstream-page}

The page has, from top to bottom: the top bar, the overview strip, the time axis, the lanes (with
the sidebar on their right), and the status bar. Before a recording is opened, the lanes show
\ui{Open a recording}, with the same two buttons as the top bar.

\begin{xltabular}{\linewidth}{@{}P{3.6cm}L@{}}
\caption{The controls of VitalStream's page, from left to right and from top to bottom.}
\label{tab:vitalstream-controls}\\
\toprule
{\tabhead Control} & {\tabhead Does} \\
\midrule
\endfirsthead
\toprule
{\tabhead Control} & {\tabhead Does} \\
\midrule
\endhead
\bottomrule
\endfoot
\ui{Open record\dots} (\key{O}) & Opens the file browser to open one record or one EDF file
  (\cref{sec:vitalstream-open}) \\
\ui{Open multiple records\dots} (\key{Shift}+\key{O}) & Opens the file browser to open a folder:
  every record in it and in its subfolders, as one session (\cref{sec:vitalstream-session}) \\
The recording's name & \emph{name} $\cdot$ \emph{format} $\cdot$ \emph{description}, for
  instance \code{session $\cdot$ WFDB records}; the recording's path on hover. Before one is
  open: \code{No recording open}. A host's context (\code{who}, \code{what}) takes its place
  when another program opens the page (\cref{sec:vitalstream-address}) \\
\ui{Go to\dots} (\key{G}) & A date and a time to go to (\cref{sec:vitalstream-goto}); available once
  a recording is open \\
The busy indicator & Shown while data are being fetched \\
\ui{Time zone} & The display time zone (\cref{sec:vitalstream-time}) \\
\ui{S} \ui{M} \ui{L} & The height of the lanes: compact, normal, tall \\
The half-filled circle & Switches between the light and the dark theme. Until it is used, the
  page follows the system's setting \\
\ui{?} (\key{?}) & The help: the keys and mouse actions, and what the shadings mean \\
The overview strip (\ui{OVERVIEW}) & The whole recording: one row per kind of signal with the
  data it holds, the alarms, the breaks, and a frame around the part in view. Click or drag
  to move the view there; the wheel zooms around the middle of the view \\
The time axis & The date and the zone's abbreviation, the spacing of the ticks
  (\code{ticks every 2 s}), the times, and the cursor's time \\
The lanes & The alarms, the waveforms and the trends (\cref{sec:vitalstream-lanes}) \\
\ui{At cursor} & The time of the cursor (else of the middle of the view), and the latest value
  of every vital sign at that time, no older than 10 minutes, with its age (\code{now}, or how
  long before); the trends shown in lanes come first. \code{No values in the 10 minutes
  before.} when there are none \\
\ui{Show all waveforms and trends} & Every waveform (up to 64) and every trend; untick it to
  return to the selection the recording opened with \\
\ui{Waveforms} & The channels, grouped by kind (\ui{ECG}, \ui{Pressure},
  \ui{SpO$_2$ / pleth}, \ui{Respiration}, \ui{CO$_2$}, \ui{Gases}, \ui{Ventilator},
  \ui{EEG, EOG, EMG}, \ui{ICG}, \ui{SvO$_2$}), each with its color, unit (\code{adu} without
  one) and sampling rate. The channels of plain WFDB records and EDF files, whose readers state no
  kind, are grouped by their labels (\code{ECG}, \code{II} and \code{V1} as ECG, \code{SpO2} and
  \code{Pleth}, \code{Airflow} and \code{Thorax} as respiration, \code{EEG}, \code{EOG} and
  \code{EMG}, \code{ART} as pressure, \code{CO2}); a channel whose label names no kind is listed
  under \ui{Signals}. Tick or untick a channel to show or hide its lane \\
\ui{Trends} & The vital-sign series, with the unit and the number of values of each; tick or
  untick a series to show or hide its lane. \ui{Filter vital signs\dots} narrows the list \\
\ui{Recording} & \ui{Format}, \ui{Start}, \ui{End}, \ui{Duration}, \ui{Waveforms} (the channels and
  the contiguous pieces they are made of), \ui{Vital signs}, \ui{Alarms} (\ui{Annotations} for an
  EDF+ file), \ui{Breaks},
  \ui{Markers} (their number and names), \ui{Time stamps} (the resolution of the source's time
  stamps, when known); then the recording's notes, and what could not be shown
  (\code{Not shown: \dots}) \\
The status bar & On the left, the span in view (start $\rightarrow$ end $\cdot$ length $\cdot$
  samples per pixel, or pixels per sample, of the fastest waveform shown); on the right, the
  messages (\code{updated in 8 ms}). Errors are written in bold: red is kept for alarms \\
\end{xltabular}

The page remembers the following in the browser: the theme (once chosen), the time zone, the
lanes' height, the last folder of the file browser, the recent folders and the recent
recordings.

% ==========================================================================================
\section{Opening a recording}
\label{sec:vitalstream-open}

Two buttons at the top of the page, repeated in the empty page, open recordings
(\cref{tab:vitalstream-open}).

\begin{table}[htbp]
\centering
\caption{VitalStream's two ways to open recordings.}
\label{tab:vitalstream-open}
\begin{tabularx}{\linewidth}{@{}P{4.4cm}L@{}}
\toprule
\tabhead Button & \tabhead Opens \\
\midrule
\ui{Open record\dots} (\key{O}) & One record: a WFDB record (its \file{.hea} file, which names its
  \file{.dat} or \file{.mat} signal files; a multi-segment record whole, or one of its segments
  alone, by the segment's own header), an EDF, EDF+, BDF or BDF+ file, or a file in a format that
  an installed adapter reads. A signal file typed as the path is refused with the header to open
  in its place (\cref{sec:recordings-one}) \\
\ui{Open multiple records\dots} (\key{Shift}+\key{O}) & A folder: every WFDB record in it and in
  its subfolders, on one time axis (one session), or a WFDB collection \\
\bottomrule
\end{tabularx}
\end{table}

\subsection{The file browser}
\label{sec:vitalstream-browser}

Both buttons open the file browser of the computer the viewer runs on (titled \ui{Open record}
or \ui{Open multiple records}), at the last folder used; the first time, at the home folder when
the viewer may open it, else at its first data root. The browser shows only what lies inside the
viewer's data roots (\cref{chap:viewers-alone}).

\begin{itemize}
\item \textbf{The path box.} Type or paste a path. \key{Tab} completes it (to the longest common
  beginning of the paths that match, and lists them), \key{$\downarrow$} and \key{$\uparrow$} pick
  one of the suggestions (up to 50), \key{Enter} goes there, \key{Esc} closes the suggestions, and
  \key{Alt}+\key{$\uparrow$} goes up one folder. \ui{Go} goes to the path typed; \ui{$\uparrow$} goes
  up one folder. A path to a record or a recording file opens it.
\item \textbf{The places} below it: \ui{Home}, \ui{sample\_data} (run from a clone), each data root
  other than the home folder and the whole file system, then the recent folders (up to six).
  Below them, the path as links, one per folder.
\item \textbf{The list.} Folders first, then records and recording files, by name; a row \ui{..}
  leads to the folder above. Folders are entered with a click; records and recording files open
  with a click. Each entry is marked:
  \begin{itemize}
  \item a folder;
  \item a folder that holds records itself, or a WFDB collection, tagged \code{WFDB records};
  \item a WFDB record (by its name, without \file{.hea}), tagged \code{WFDB record} or
    \code{WFDB multi-segment record};
  \item an EDF or BDF file, tagged with its format (\code{EDF}) and its size;
  \item a file an adapter reads, tagged with its format's label and its size;
  \item a symbolic link, with \code{link}; a link that leads nowhere, tagged \code{broken link}.
  \end{itemize}
  Signal files, files no installed reader reads, and the segments and layout headers of the
  multi-segment records in the folder are not listed. Nothing outside the data roots is listed,
  symbolic links that lead outside among them. A folder of more than 5000 entries shows its first
  5000 (\code{Showing the first 5000 entries; type a path or use the filter to narrow down.}).
\item \ui{Filter this folder\dots} narrows the list to the names that contain what is typed;
  \key{Enter} opens the only one left.
\item \ui{hidden}: also list the hidden entries (names that start with \code{.}) and the segments
  of multi-segment records.
\item \ui{Recent recordings}: the last eight recordings opened, by path; a click opens one again.
\end{itemize}

A path that cannot be shown is answered in the browser: \code{path outside the configured data
roots}, \code{no such file or folder: PATH}, \code{not a folder or a recording this viewer reads:
PATH}, \code{permission denied: PATH}.

\subsection{Open multiple records: the folder's summary}
\label{sec:vitalstream-summary}

In \ui{Open multiple records\dots}, each folder of the list has an \ui{Open} button, and
\ui{Open the records in this folder} opens the folder being shown. Before anything is read, the
viewer counts the records (\code{Counting the records\dots}) and shows a summary of the folder,
from its headers alone (\cref{fig:vitalstream-summary}):
\begin{itemize}
\item the records, the segments of multi-segment records, and the signals (their names, the first
  twelve);
\item the time span (\code{From \dots{} to \dots}), and the length of data in all; or, when no record
  states its time, the length of data alone;
\item the records that state no time (\code{placed after the data before});
\item the records that cannot be read (\code{will be left out});
\item the size to index (\code{to index once; later openings are instant.}).
\end{itemize}
\ui{Open as one session} opens it; \ui{Back} returns to the browser. A folder over a limit, or
without records, is refused at this step (\cref{sec:vitalstream-limits}).

\screenshot{vitalstream-open-multiple.png}{The summary of a folder, before it is opened as one session: six records, four segments of multi-segment records, three signals, the time span, two records that state no time, and the size to index.}{fig:vitalstream-summary}

\subsection{Indexing, and Cancel}
\label{sec:vitalstream-indexing}

The first opening of a recording indexes it: a progress bar shows how far it is (\code{Indexing
NAME \dots{} 40 \%{} $\cdot$ about 12 s left. This is done once; the index is cached.}), and
\ui{Cancel} stops it: \code{Cancelled: nothing was kept.} Later openings are instant: the index
is kept in the index folder (\cref{sec:vitalstream-index}). Once it is ready, the dialog closes
and the recording opens, whole, in the lanes. A recording that cannot be read is reported in the
dialog (\code{Error: \dots}) and in the status bar (\code{Could not open: \dots}).

\subsection{A folder as one session}
\label{sec:vitalstream-session}

A folder without an \file{index.csv} is read as one session: every record in it and in its
subfolders, on one time axis. This is how a patient's records, split into folders by device or by
day, are seen together. A folder with an \file{index.csv} is a WFDB collection, read as its files
describe it (\cref{chap:formats-reference}).

\begin{itemize}
\item \textbf{Records} are the \file{.hea} files below the folder, at every level, in the order of
  their path relative to the folder (\code{/} between folders, compared character by character).
  Hidden entries (names that start with \code{.}) and symbolic links are skipped; the segments and
  the layout header of a multi-segment record belong to that record, never records of their own.
  A \file{RECORDS} file is not needed, and not read; a \file{vitals.csv}, \file{events.csv} or
  \file{gaps.csv} in the folder itself is read as for a collection.
\item \textbf{Placement.} A record starts at its header's base date and time, taken as UTC. A
  record without one is placed after the data before it, and shaded \emph{time not stated}
  (\cref{sec:vitalstream-untimed}). One record alone without a time starts at 0
  (1970-01-01).
\item A record that cannot be read, or whose header states no sampling frequency, is left out;
  the recording's notes say how many.
\end{itemize}

\paragraph{Multi-segment records.} A multi-segment record is read as one record. Its segments
follow each other from the record's start, each at its offset (the samples of the segments
before it, divided by the sampling frequency); each signal of each segment is a piece of the
channel of its name. A null segment (\code{\textasciitilde}) is a break of known length, with
the reason \code{null segment}. In a record with a variable layout, the layout segment (the first,
of length 0) names the signals, and a signal absent from a segment is a break in its channel
only. The record's markers are those of the master header's \file{.mrk} file; a marker in a null
segment is left out. Opening one segment's own header (\file{rec\_0001.hea}) reads that segment
alone, as one record.

\paragraph{Overlapping records.} A signal recorded twice at the same time (two records of lead II
whose times overlap) is never drawn on top of itself: each piece goes into the first lane of its
channel that it does not overlap, and each further lane is a channel of its own, named with its
number: \code{II}, \code{II (2)}, \code{II (3)}. Pieces are placed in the order of the records,
and the signals of a record in the order of its header. The recording's notes say so when it
happens. A channel is a kind of signal (\code{wfdb} for plain records), a signal description and
a unit.

\screenshot{vitalstream-session.png}{The synthetic folder \file{testdata/session} opened as one session: the lanes ART, II, II (2) and V; a span whose time is not stated (amber); the break of a null segment (grey); R-wave markers.}{fig:vitalstream-session}

\subsection{The limits on a folder}
\label{sec:vitalstream-limits}

So that opening a large folder by mistake cannot overwhelm the viewer, a folder of plain records
is opened only within these limits. They are VitalStream settings (\code{max\_records},
\code{max\_depth}, \code{max\_scan\_s}, \cref{sec:vitalstream-settings}); Medorama has no such
settings, and indexes a recording's folder within the defaults:

\begin{table}[htbp]
\centering
\caption{The limits on a folder opened as one session.}
\label{tab:vitalstream-limits}
\begin{tabularx}{\linewidth}{@{}LP{1.6cm}P{5.4cm}@{}}
\toprule
\tabhead Limit & \tabhead Default & \tabhead Over it \\
\midrule
Header files below the folder (segments and layout headers counted) & 5000 &
  \code{the folder holds more than 5000 WFDB records} \\
Folder levels below the folder & 20 & \code{the folder has records nested more than 20 levels
  deep} \\
Seconds to count the records (the summary) & 10 & \code{counting the records took more than
  10 s} \\
\bottomrule
\end{tabularx}
\end{table}

Over a limit, the viewer refuses the folder, says which limit was reached, and what to do
instead: \code{open a subfolder, or describe the records in an index.csv (a WFDB collection,
docs/wfdb-collection.md)}. A WFDB collection lists its records itself and is not limited by
these. A folder without any record is refused too (\code{no WFDB records in or below PATH}).
Every recording, a collection included, has at most 100,000 pieces (signals $\times$ records,
or the runs of an EDF+D file), and a folder's index key is computed from at most 200,000 of its
files within \code{max\_depth} levels (\cref{sec:formats-reference-limits}).

\subsection{Trying it}
\label{sec:vitalstream-try}

The synthetic recordings of \file{viewers/vitalstream/testdata/} (no patient data) show what the
viewer does: \ui{Open multiple records\dots}, then \file{viewers/vitalstream/testdata/session}
(\cref{fig:vitalstream-session}), \file{testdata/collection} (a WFDB collection with vital signs,
alarms and breaks) or \file{testdata/plain} (plain records). The clone must lie inside the viewer's
root (by default, the home folder).

% ==========================================================================================
\section{Moving through a recording}
\label{sec:vitalstream-navigate}

A recording opens whole. The mouse and the keys move and zoom the view; the keys work when the
focus is not in a text box, and are ignored with \key{Ctrl}, \key{Alt} or \key{Cmd}.

\begin{xltabular}{\linewidth}{@{}P{5.2cm}L@{}}
\caption{The mouse actions and keys of VitalStream.}
\label{tab:vitalstream-keys}\\
\toprule
{\tabhead Action} & {\tabhead Effect} \\
\midrule
\endfirsthead
\toprule
{\tabhead Action} & {\tabhead Effect} \\
\midrule
\endhead
\bottomrule
\endfoot
Wheel, or a trackpad pinch, over the traces & Zoom at the pointer \\
Horizontal scroll over the traces & Pan \\
Drag & Pan \\
\key{Shift} + drag & Zoom to the time span selected \\
Double-click, \key{Home} & The whole recording \\
Click or drag in the overview strip & Move the view there \\
Wheel over the overview strip & Zoom around the middle of the view \\
Wheel over the lane labels & Scroll the lanes \\
\key{$\leftarrow$} \key{$\rightarrow$} & Pan 10\,\% of the view \\
\key{PgUp} \key{PgDn} & Pan one screen \\
\key{+} (or \key{=}) \key{$-$} & Zoom in and out, by a factor of 2, around the cursor (else the
  middle of the view) \\
\key{1} \key{2} \key{3} & Show 10\,s, 1\,min or 10\,min around the cursor (else the middle of the
  view) \\
\ui{Go to\dots}, \key{G} & A calendar of the recording's days and a time, to the millisecond
  (\cref{sec:vitalstream-goto}) \\
\key{O} $\cdot$ \key{Shift}+\key{O} $\cdot$ \key{?} & Open a record $\cdot$ open multiple records
  $\cdot$ help \\
In the file browser: \key{Tab} $\cdot$ \key{Enter} $\cdot$ \key{Alt}+\key{$\uparrow$} & Complete
  the path $\cdot$ go there $\cdot$ go up one folder \\
\end{xltabular}

The view can be zoomed in to 20\,ms, and out to a little more than the whole recording; it can
be moved a little beyond the data on each side. The pointer over the traces is the cursor: a
vertical line, with its time on the axis, the value of each lane at the right of the lane, and
the vital signs at that time in \ui{At cursor}.

\subsection{Go to a date and a time}
\label{sec:vitalstream-goto}

\ui{Go to\dots} (\key{G}) opens a calendar at the cursor's time (else the middle of the view):
\begin{itemize}
\item the month, with \ui{\textless} and \ui{\textgreater} for the months before and after (within
  the recording); weeks start on Monday; the days with waveforms are marked, the days outside the
  recording cannot be chosen;
\item \ui{Time}, to the millisecond, in the display time zone (named beside it);
\item \ui{Show}: \ui{the current span}, \ui{10 s}, \ui{1 min}, \ui{10 min} or \ui{1 h}
  \emph{around it};
\item \ui{Start} and \ui{End}: the recording's first and last moment;
\item \ui{Go} (or \key{Enter}, or a double click on a day) goes there; \key{Esc} closes it.
\end{itemize}
The cursor line then marks the time asked for, and the status bar says \code{Went to \dots}. A
time outside the recording is refused, with the recording's span.

% ==========================================================================================
\section{The lanes}
\label{sec:vitalstream-lanes}

The lanes are, from top to bottom: the alarms (when the recording has any), the waveforms,
grouped by kind (ECG, pressures, SpO$_2$ and pleth, respiration, CO$_2$, gases, ventilator, EEG,
ICG, SvO$_2$, then the others), and the trends. For plain WFDB records and EDF files, the kind of
each channel is recognized from its label; this affects only the grouping and the colors on the
page, not the index. Each lane's label gives its name, its unit and,
for a waveform, its sampling rate, with the top and bottom of its scale. Values are in the
units the file states; \code{adu} means it states no calibration.

\subsection{Waveforms, and the 64-lane cap}
\label{sec:vitalstream-waves}

Waveform lanes show, for long views, the minimum and the maximum of the samples behind each
pixel column, so that no peak is lost, and the samples themselves when zoomed in (with a dot on
each sample once they are far enough apart). The y-axis fits the data in view. A missing sample
is a gap in the line, never a zero. A lane without data in view says \code{no data in view}.

At most 64 waveform lanes show at once. A recording with more opens with the first 64, in lane
order, and the status bar says so (\code{N waveforms: the first 64 are shown; pick others under
Waveforms.}); another is shown by unticking one under \ui{Waveforms} and ticking it. Ticking a
65th is refused (\code{At most 64 waveform lanes at once: untick one first.}).

\subsection{Trends and vital signs}
\label{sec:vitalstream-trends}

Trend lanes show the vital signs as stored, as points joined by a line (not across more than 15
minutes without a value); when there are more values than pixels, the minimum and maximum of
each pixel column. The value under the cursor is shown only when a stored value lies within the
15 minutes before it. A recording opens with these trends, when it has them: heart rate (\code{HR}),
SpO$_2$, respiration rate, the mean arterial pressure and the mean non-invasive pressure;
\ui{Show all waveforms and trends}, or the boxes under \ui{Trends}, show the others. A label
measured by several sources is a series per source, never merged: \code{HR}, \code{HR [import]}.

\subsection{Alarms}
\label{sec:vitalstream-alarms}

Alarms are bands in the lane \ui{Alarms}, colored by severity: 1 advisory, 2 warning, 3 crisis.
They are the only red on the page. The annotations of an EDF+ file are shown in the same way, in
a lane named \ui{Annotations}; they have no severity. The monitor's alarm state, when the recording has it, has a row
of its own above the other alarms. A band shows its text when it is wide enough; on hover, its
text, its kind and level, its start and end, and its length. The overview strip shows the alarms
too.

\subsection{Markers}
\label{sec:vitalstream-markers}

Markers, such as R waves and pacing spikes, are ticks at the top of their channel's lane, in
views up to 2 minutes. With fewer than 80 in a lane, each is labelled: \code{R} for an R wave,
\code{P}, \code{pace}, \code{defib}, or the marker's name.

\subsection{Breaks}
\label{sec:vitalstream-breaks}

A break is a place where contiguity is not proven. Breaks are shaded grey: across every lane
for a break of the whole recording, only across a channel's lanes for a break of that channel. On
hover: \code{Break in the data:} and its reason, its length (\code{(exact)} when it is) and how it
was measured, or \code{gap duration unknown}, its start and end, and its notes. Pieces of a
channel are drawn apart, never joined across a break.

\subsection{Time not stated}
\label{sec:vitalstream-untimed}

Data whose time the recording does not state is shaded amber, labelled \code{TIME NOT STATED}
where it is wide enough. It is placed after the data before it; its place is known relative to
its neighbors, not on the clock, and the timing of its samples within it is exact. On hover:
\code{Time not stated in the file for this data.}

\subsection{Repaired data: \texorpdfstring{\code{qc\_verbose}}{qc\_verbose}}
\label{sec:vitalstream-qc}

Some readers repair what a file stores out of order: the STP adapter puts back in order the data
blocks a file stored late, so that they make no break. By default, such repaired data is not
marked. With the setting \code{qc\_verbose} (or \opt{-{}-qc-verbose}), it is outlined in dashed
blue, with its details on hover (\code{Data block stored late in the file (N s, N blocks); the
reader put it back in order. Not a break: shown because qc\_verbose is on.}). Breaks are shaded
either way.

% ==========================================================================================
\section{Time zones}
\label{sec:vitalstream-time}

Recordings are kept in UTC, as Unix seconds; a WFDB header's base date and time is taken as UTC.
The time zone only changes how times are shown: on the axis, at the cursor, in \ui{Recording},
in \ui{Go to\dots} and in the address's view.

The list \ui{Time zone}, at the top right, holds the viewer's default zone, marked
\code{(default)}, then \code{UTC}, then every zone the browser knows. The default is the setting
\code{timezone} (or \opt{-{}-timezone}), \code{America/New\_York} unless set; a choice in the list
is remembered in the browser. The axis names the zone's abbreviation at the time shown
(\code{EST}, \code{EDT}).

% ==========================================================================================
\section{The address}
\label{sec:vitalstream-address}

The page's address keeps what it shows in its fragment (after \code{\#}), so that a reload or a
bookmark returns to the same place, and another program can open the viewer at a recording:

\begin{table}[htbp]
\centering
\caption{The fragments of VitalStream's address.}
\label{tab:vitalstream-address}
\begin{tabularx}{\linewidth}{@{}P{5.0cm}L@{}}
\toprule
\tabhead Fragment & \tabhead Means \\
\midrule
\code{path=\dots\&t0=\dots\&t1=\dots} & A recording on the server (its path), at a time span (Unix
  seconds) \\
\code{browse=FOLDER} & The file browser, open at a folder \\
\code{who=\dots}, \code{what=\dots}, \code{pos=\dots}, \code{nav=\dots} & A host's context, shown in
  place of the recording's name (\cref{chap:viewers-alone}) \\
\code{theme=light} or \code{dark} & The host's colors: applied, not saved as the viewer's own
  choice (the theme button still switches them); Medorama passes its page's theme this way
  (\cref{chap:viewers-in-medorama}) \\
\bottomrule
\end{tabularx}
\end{table}

The page's requests to its server are described in \cref{chap:viewer-api}.

% ==========================================================================================
\section{Settings}
\label{sec:vitalstream-settings}

A settings file, \file{vitalstream.toml}, changes the defaults. It is \opt{-{}-config FILE},
else the first that exists of \file{viewers/vitalstream/vitalstream.toml} (only for a viewer run
from the build folder of a clone: git ignores it) and
\file{\textasciitilde/.config/vitalstream/vitalstream.toml}; only that one is read. A
\file{vitalstream.toml} in the working folder is never read unless \opt{-{}-config} names it: a
shared data folder could hold one that runs programs (\code{adapters\_dir}) or opens the viewer
to the network (\code{host}, \code{auth\_file}). On start, the viewer says which settings file it
read. One file serves the three viewers: the top-level keys apply to all, the keys
of a \code{[py]}, \code{[rs]} or \code{[cpp]} section to that viewer only, and override the
top-level ones. Command-line options override both. Every key is optional; relative paths are
relative to the file's folder, and \code{\textasciitilde} is the home folder. A key the viewer does
not know is an error (\code{unknown settings: KEY}), and so is a value of the wrong type
(\code{bad value for KEY}): the viewer does not start. \code{data\_roots} and
\code{allowed\_hosts} are lists of strings, \code{port} an integer from 0 to 65535,
\code{index\_workers}, \code{max\_records} and \code{max\_depth} integers of at least 0,
\code{max\_scan\_s} a number of at least 0, the switches \code{true} or \code{false}, the others
strings. Start from
\file{viewers/vitalstream/vitalstream.example.toml}.

\begin{xltabular}{\linewidth}{@{}P{3.4cm}P{3.0cm}L@{}}
\caption{The settings of VitalStream (\file{vitalstream.toml}).}
\label{tab:vitalstream-settings}\\
\toprule
{\tabhead Key (option)} & {\tabhead Default} & {\tabhead Meaning} \\
\midrule
\endfirsthead
\toprule
{\tabhead Key (option)} & {\tabhead Default} & {\tabhead Meaning} \\
\midrule
\endhead
\bottomrule
\endfoot
\code{data\_roots} (\opt{-{}-root}) & The viewer's root & The folders the viewer may open, inside
  its root: one that contains the root is narrowed to it, one outside it is left out (the
  settings narrow the root, never widen it). A path is allowed when the file or folder it leads
  to, every symbolic link resolved, lies inside a data root; a link that leads outside is
  refused, and not listed \\
\code{max\_records}, \code{max\_depth}, \code{max\_scan\_s} & 5000, 20, 10 & The limits on a
  folder of records: header files, folder levels, seconds to count them
  (\cref{sec:vitalstream-limits}); \code{max\_scan\_s} is a number of seconds, 0 or more \\
\code{require\_password} & \code{false} & Refuse to start without a password (Medorama sets it
  for the viewer it runs; a viewer built without password support then never starts) \\
\code{timezone} (\opt{-{}-timezone}) & \code{America/New\_York} & The display time zone, an IANA
  name; data are stored in UTC \\
\code{host}, \code{port} (\opt{-{}-host}, \opt{-{}-port}) & \code{127.0.0.1}; 8050, 8060, 8070 &
  Where the viewer listens \\
\code{allowed\_hosts} & \code{[]} & The host names the viewer answers for besides
  \code{localhost}, \code{127.0.0.1}, \code{::1} and \code{host}: the names this machine is
  reached by when \code{host} is \code{0.0.0.0}, or a reverse proxy's public name;
  \code{["*"]} allows any. A request naming another host is refused, which defeats DNS
  rebinding (\cref{sec:viewer-api-checks}) \\
\code{secure\_cookie} & \code{false} & Marks the sign-in cookie \code{Secure}, so that the
  browser sends it over HTTPS only: set it for a viewer behind an HTTPS proxy \\
\code{index\_workers} & 2 & The recordings indexed at the same time (0 is taken as 1) \\
\code{cache\_dir} (\opt{-{}-cache-dir}) & \file{\textasciitilde/.cache/vitalstream} (run from a
  clone: \file{viewers/vitalstream/sample\_data/outputs/index}) & The index folder; shared by the
  three viewers, and with Medorama when set to its \code{index\_dir} \\
\code{adapters\_dir} (\opt{-{}-adapters-dir}) & \file{\textasciitilde/.local/share/medorama/adapters}
  (Windows: \code{\%LOCALAPPDATA\%\textbackslash medorama\textbackslash adapters}) & The adapter
  folder \\
\code{auth\_file} & \file{\textasciitilde/.config/vitalstream/auth.json} & The password: a salted
  scrypt hash, readable by you only \\
\code{qc\_verbose} (\opt{-{}-qc-verbose}) & \code{false} & \code{false}: data the reader repaired
  (blocks stored out of order and put back in order) is not marked; \code{true} outlines it
  (dashed), with its details on hover. Breaks are shaded either way \\
\end{xltabular}

An example, for a server whose recordings lie in two folders:

\begin{lstlisting}[language=TOML,style=example]
# ~/.config/vitalstream/vitalstream.toml, for a viewer installed with --root /data
timezone = "Europe/Paris"
data_roots = ["/data/icu", "/data/holter"]
host = "127.0.0.1"
index_workers = 4
max_records = 10000

[cpp]
port = 8070
\end{lstlisting}

The settings of every program of Medorama, VitalStream's among them, are gathered in
\cref{chap:settings-reference}.

% ==========================================================================================
\section{The index}
\label{sec:vitalstream-index}

The index of a recording is what lets the viewer show days of waveforms at once: the
recording's structure (its channels, and the pieces each is made of), the minimum and maximum
envelopes of every waveform, and its markers, vital signs, alarms and breaks. A recording is
indexed the first time it is opened, or ahead of time with \opt{-{}-index PATH}; later openings
are instant.

\begin{itemize}
\item The index folder (\code{cache\_dir}) holds one folder per recording, named by its
  \emph{key}: the first 20 hexadecimal digits of a SHA-256 of the index format's version, the
  recording's real path, and its files' sizes and modification times. A recording that is changed, moved or replaced gets a new
  key, so an outdated index is never used.
\item The three viewers and Medorama write the same index for the same recording, byte for byte,
  and may share one index folder at the same time.
\item An index is written into a temporary folder and renamed into place, so that a reader sees
  a complete index or none; it is never written next to a recording.
\end{itemize}

\begin{phibox}
An index holds no name, identifier or date of birth, but it holds the recording's path and its
values, which are patient data. Keep index folders on storage protected like the recordings.
\end{phibox}

The index format is specified in \cref{chap:formats-reference}.

% ==========================================================================================
\section{Deploying}
\label{sec:vitalstream-deploy}

\paragraph{On your computer:} install, then run the command; the page opens in your browser. On a
machine others can log in to, set a password once (\opt{-{}-set-password}): anyone who can connect
to the port can otherwise read the files under the viewer's root.

\paragraph{On a server where the data are:} the data need not be copied anywhere.
\begin{enumerate}
\item Install it with a root that holds the recordings, and set a password once:
  \cmd{viewers/vitalstream/install.sh -{}-root /data/icu}, then
  \cmd{vitalstream-cpp -{}-set-password}.
\item Start it without a browser, in \cmd{tmux} or as a service (below):
  \cmd{vitalstream-cpp -{}-no-browser}.
\item From your computer: \cmd{ssh -N -L 8070:127.0.0.1:8070 user@server}, then open
  \url{http://127.0.0.1:8070}.
\end{enumerate}

The viewers listen on 127.0.0.1 by default, keep a sign-in for 30 days, and work behind a portal
or a reverse proxy that serves them under a path prefix (their page uses relative addresses).
When the sign-in ends, or the password changes, the page shows the sign-in form again. Signing
out ends that sign-in on the server too, so that a copy of the cookie no longer works. After
wrong passwords from one address, the next one is checked only after a delay (1, 3, 7, 15
\dots{} up to 30\,s); the C++ viewer answers such an early password \code{429}
(\cref{sec:viewer-api-signin}). A viewer that answers for other names than this machine's local
ones (\code{host = "0.0.0.0"}, or a reverse proxy that passes its public name on) lists them in
\code{allowed\_hosts}; behind an HTTPS proxy, set \code{secure\_cookie = true}.

\paragraph{Run as a service} (Linux, systemd, as your user):

\begin{lstlisting}[style=shell]
mkdir -p ~/.config/systemd/user
cat > ~/.config/systemd/user/vitalstream.service <<'EOF'
[Unit]
Description=VitalStream viewer

[Service]
ExecStart=%h/.local/bin/vitalstream-cpp --no-browser
Restart=on-failure

[Install]
WantedBy=default.target
EOF
systemctl --user daemon-reload && systemctl --user enable --now vitalstream
loginctl enable-linger "$USER"   # keep it running after you log out (may need an administrator)
\end{lstlisting}

On macOS and Windows, a viewer runs as a service as Medorama does, with the viewer's command
(\cref{chap:deployment}).

% ==========================================================================================
\section{Adapters}
\label{sec:vitalstream-adapters}

VitalStream reads WFDB records and EDF files itself. Every other recording format is read by an
\emph{adapter}: a separate program, named \file{NAME-adapter} (\file{NAME-adapter.exe} on
Windows), in the adapter folder (the setting \code{adapters\_dir}, by default
\file{\textasciitilde/.local/share/medorama/adapters}, the folder Medorama uses too). Formats that are
licensed, such as the \file{.Stp} files of BedMasterEx, are distributed this way, as add-ons,
outside the repository; the stp-adapter that reads them is private and available to institutions
under agreement. The built-in readers take precedence over an adapter that lists the same
extensions.

\begin{itemize}
\item \textbf{At start}, the viewer asks every adapter of its folder to describe itself, once:
  the formats it reads, each with a kind, a label and file extensions.
\item \textbf{Browsing.} A file whose extension an adapter lists is a recording of that adapter's
  kind, listed with its format's label and size. Until its adapter is installed, a file of that
  format is treated like any other file the viewer does not read: it is not listed, and opening
  it is refused.
\item \textbf{Indexing.} The viewer runs the adapter, which writes the index into a temporary
  folder and reports its progress. The viewer then completes the index and renames the folder
  into place.
\item \textbf{Viewing.} Envelopes, markers and vital signs come from the index. Raw samples, when
  a view zooms in close enough, are read from the recording through the adapter.
\end{itemize}

To install an add-on for a viewer alone, copy its program into the adapter folder, make it
executable on Linux and macOS (\cmd{chmod 755}), allow the system to run it (on macOS,
\cmd{xattr -d com.apple.quarantine PROGRAM}; on Windows, \cmd{Unblock-File PROGRAM}), and restart
the viewer. Installing add-ons for Medorama, and the adapters' protocol, are described in
\cref{chap:addons}.

% ==========================================================================================
\section{Testing and parity}
\label{sec:vitalstream-testing}

\begin{xltabular}{\linewidth}{@{}P{5.6cm}L@{}}
\caption{VitalStream's tests (from the root of the repository).}
\label{tab:vitalstream-testing}\\
\toprule
{\tabhead Command} & {\tabhead Checks} \\
\midrule
\endfirsthead
\toprule
{\tabhead Command} & {\tabhead Checks} \\
\midrule
\endhead
\bottomrule
\endfoot
\cmd{cargo test -p vitalstream-rs -p recording-engine} & The Rust viewer and engine on the
  synthetic recordings of \file{testdata/} (one record, a folder read as one session, the
  limits, a record's index key, cancelling, headers that state too many samples, files outside
  the data roots); the viewer's root and its settings files; the HTTP API's checks and sign-in;
  the adapter protocol with a mock adapter and misbehaving ones \\
\cmd{(cd viewers/vitalstream/cpp \&\& ctest -{}-preset release)} & The C++ viewer and engine: the
  same recordings and adapter; the engine's limits and checks (\file{test\_limits.cpp}) and the
  HTTP server's (\file{test\_server.cpp}) \\
\cmd{pytest viewers/vitalstream/python/vitalstream/tests} & The Python viewer: the same recordings
  and adapter, sign-in, the file browser, a build without password support, and the limits and
  checks shared with the Rust engine and viewer (\file{test\_limits.py}) \\
\cmd{make parity} & The viewers that are built: byte-identical indexes under the specified key,
  and the same API answers (\file{viewers/vitalstream/tests/parity}) \\
\end{xltabular}

\cmd{make parity} runs \file{tests/parity/index\_parity.py}, which indexes recordings with every
viewer that is built and checks that each key equals the one computed from the specification,
by an independent implementation, and that the index files are byte-identical (\file{meta.json}
compared without its informative fields), and \file{tests/parity/viewer\_parity.py}, which checks
that the viewers give the same answers to the same requests, and the same statuses and errors
to the requests they refuse (the \code{Host} check, \code{Origin}, \code{415}, \code{413}, the
query grammar, a record whose signal file leads outside the data roots). The first also takes recordings an
adapter reads:

\begin{lstlisting}[style=shell]
python viewers/vitalstream/tests/parity/index_parity.py                            # testdata/
python viewers/vitalstream/tests/parity/index_parity.py --adapters-dir DIR FILE... # an adapter's
\end{lstlisting}

\file{testdata/generate.py} writes the synthetic recordings, deterministically, with the Python
standard library only: a WFDB collection with three stretches, a break of stated length, a record
without a start time, vital signs from two sources, alarms and markers; plain WFDB records, among
them an uncalibrated signal and a record without a sampling frequency; a folder of records in
subfolders to open as one session (two overlapping ECG records, a record without a start time,
multi-segment records with a null segment and a variable layout, a hidden folder that is not
read); and four EDF files in \file{testdata/edf/} (\cref{sec:formats-reference-edf-tests}).

% ==========================================================================================
\section{License}
\label{sec:vitalstream-licence}

VitalStream is licensed, like all of Medorama, under the PolyForm Noncommercial License 1.0.0,
with commercial licenses by agreement. The three viewers read WFDB records with wfdb-x, a
separate project whose library is under the GNU LGPL version 2.1 or later (\cref{chap:licences}).
